% frontmatter/slowa/verano.tex -- la página para el adulto del cuaderno
% de verano de «Pierwsze słowa» (ediciones/slowa-verano.tex, ver
% notes/07-ediciones.md): va en vez de frontmatter/slowa/como-usar.tex,
% que habla del año entero. Detrás, en vez del mapa del año, el mapa del
% verano (\mapaEdicion, preamble.tex).
\section*{Jak korzystać z tego zeszytu na lato}

To lato z \emph{Pierwszych słów}, pierwszego zeszytu z serii
\emph{Uczę się czytać} po polsku: jego ostatnie 13 tygodni, w osobnym
zeszycie. To lato, w którym od czytania słów przechodzi się do czytania
zdań: codziennie jedno krótkie zdanie, od dwóch do czterech słów
(\emph{Toby czeka.}, \emph{Bigotes śpi.}). Jest dla dziecka, które
czyta już słowa z kilkoma sylabami -- także z dwuznakami (\emph{sz, cz,
rz, ch, dz}), z \emph{ą} i \emph{ę} i z dwiema spółgłoskami obok siebie
(\emph{szko·ła}, \emph{kro·wa}) -- i zaczyna łączyć je w zdania. Na
końcu zeszytu będzie tam, gdzie zaczyna się zeszyt ze zdaniami.

Na każdy dzień od poniedziałku do piątku jest jedna strona, razem 65.
Wystarczy pięć albo dziesięć minut dziennie. Każda strona wygląda tak
samo:

\begin{itemize}[leftmargin=6mm]
\item \textbf{Data}, \textbf{Czytanie} i \textbf{Uwagi}: dla dorosłego.
  Zaznaczenie codziennie, czy dziecko czytało samodzielnie, z pomocą czy
  z trudem, nic nie kosztuje, a pod koniec lata jest zapisem jego
  postępów.
\item \textbf{Zdanie dnia}, w zielonej ramce, dużymi literami: czyta się
  je, pokazując palcem każde słowo.
\item \textbf{Zadanie}, które prawie zawsze kończy się rysowaniem.
  Polecenia (małą szarą czcionką) \textbf{czyta dorosły} na głos:
  dziecko czyta tylko to, co jest napisane dużymi literami.
\end{itemize}

\subsection*{Zadania}

{\small
\begin{itemize}[leftmargin=6mm, itemsep=1pt, topsep=2pt]
\item \textbf{\lblDibuja}: narysować to, co się właśnie przeczytało. To
  najprostszy sposób, żeby sprawdzić, że zdanie zostało zrozumiane, a
  nie tylko odczytane.
\item \textbf{\lblCompleta}: kilka linii bez kształtu, z których trzeba
  zrobić rysunek.
\item \textbf{\lblTraza}: poprowadzić palcem albo ołówkiem po literze,
  wielkiej i małej, a potem napisać ją samodzielnie na linii.
  \textbf{\lblEscribe}: to samo z całym słowem, w środku jego pustych
  liter.
\item \textbf{\lblEncuentra}: zakreślić słowo za każdym razem, kiedy
  pojawia się między innymi, bardzo do niego podobnymi. Trzeba patrzeć
  na wszystkie litery, nie tylko na pierwszą.
\item \textbf{\lblAdivina}: dorosły czyta zagadkę, a dziecko rysuje
  odpowiedź i pisze pod spodem, co to jest.
\item \textbf{\lblUne}: połączyć linią każde zwierzę z jego odgłosem,
  każde słowo z przeciwnym albo z tym, które się z nim rymuje.
\item \textbf{\lblSiNo}: czytać bardzo krótkie zdania i decydować, czy
  mają sens (\emph{Toby szczeka} -- tak; \emph{Toby lata} -- nie).
\item \textbf{\lblRelee} i \textbf{\lblRepasa} (w piątki): wszystkie
  zdania z tygodnia razem, żeby przeczytać je jeszcze raz i zaznaczyć
  te, które wychodzą już same. Co dziesięć stron jest mały napis, żeby
  to uczcić.
\end{itemize}}

\subsection*{Jak pomagać}

{\small
\begin{itemize}[leftmargin=6mm, itemsep=1pt, topsep=2pt]
\item Jeśli dziecko utknie na jakimś słowie, warto dać mu kilka sekund,
  a jeśli nie wychodzi, przeczytać je, pokazując palcem sylaby -- i
  poprosić, żeby je powtórzyło. Lepsze to niż zamieniać stronę w walkę.
\item Nie trzeba poprawiać każdego błędu. Jeśli dziecko przeczytało
  \emph{kot} zamiast \emph{kto}, wystarczy zapytać: „na pewno? popatrz
  jeszcze raz”.
\item Zadania nie trzeba kończyć tego samego dnia. Najważniejsze jest
  czytanie; rysunek może poczekać do popołudnia.
\end{itemize}}

Strony opowiadają lato \textbf{Lucíi}, \textbf{Daniego} i ich psa
\textbf{Toby'ego} na wsi u babci Rosy, z tymi samymi tygodniami, co
hiszpański zeszyt na lato: jeśli w domu jest rodzeństwo, każde może
czytać swój zeszyt, o tym samym. Na następnej stronie są wszystkie, ze
słońcem do pokolorowania za każdy dzień czytania; na końcu zeszytu są
odpowiedzi i dyplom.
\newpage
